\documentclass{beamer}
\usepackage{amsfonts,amsmath}
\usetheme{nankai}

\usefonttheme[onlymath]{serif}
\titlebackground*{assets/background}

\title{Capital Asset Pricing Model with Size Factor}
\subtitle{Normalizing by Volatility Index}
\author{Abraham Atsiwo, Andrey Sarantsev}
\date{November 2024}

\begin{document}

\begin{frame}
\titlepage
\end{frame}

\begin{frame}{Outline}
\tableofcontents
\end{frame}

\section{Introduction}

\begin{frame}{Research Motivation}
\begin{itemize}
    \item \textbf{CAPM}: Classic model relating portfolio returns to benchmark
    \item \textbf{Size Effect}: Small stocks have higher risk and return than large stocks
    \item \textbf{Volatility Index}: Normalizing by VIX makes returns independent and normal
    \item \textbf{This Paper}: Combines all three ideas into a unified discrete-time model
\end{itemize}
\end{frame}

\begin{frame}{Key Concepts}
\begin{block}{Price Returns vs Total Returns}
\begin{itemize}
    \item \textbf{Price returns}: Only measure price changes (ignore dividends)
    \item \textbf{Total returns}: Include both price changes and dividends
    \item \textbf{Equity premium}: Total returns minus risk-free returns
\end{itemize}
\end{block}

\begin{block}{Beta and Alpha}
\begin{itemize}
    \item $\beta$: Systematic risk measure (market exposure)
    \item $\alpha$: Excess return after adjusting for market exposure
\end{itemize}
\end{block}
\end{frame}

\section{Background and Motivation}

\begin{frame}{Capital Asset Pricing Model (CAPM)}
\begin{block}{Basic CAPM Formula}
$$Q = (1-\beta)r + \beta Q_0$$
\end{block}

\begin{block}{Equity Premium Form}
$$P = \beta P_0$$
where $P = Q - r$ (equity premium) and $P_0 = Q_0 - r$ (benchmark premium)
\end{block}

\begin{itemize}
    \item $\beta = 0$: Risk-free portfolio
    \item $\beta = 1$: Market portfolio (benchmark)
    \item $\beta \in (0,1)$: Mix of bonds and stocks
\end{itemize}
\end{frame}

\begin{frame}{CAPM with Alpha}
\begin{block}{Linear Regression Model}
$$P = \alpha + \beta P_0 + \varepsilon, \quad \mathbb{E}\varepsilon = 0$$
\end{block}

\begin{itemize}
    \item $\alpha$: Excess return (alpha) - manager's skill
    \item $\beta$: Market exposure
    \item $\varepsilon$: Idiosyncratic error term
\end{itemize}

\vspace{0.3cm}
\textbf{Criticism}: CAPM assumes $\beta$ is the only risk factor, but research shows other factors matter (size, value, etc.)
\end{frame}

\begin{frame}{Size and Value Effects}
\begin{itemize}
    \item \textbf{Size factor}: Average market capitalization of portfolio stocks
    \item \textbf{Value factor}: Fundamentals (earnings, dividends, book value) vs price
    \item Well-accepted by financial community and industry practitioners
    \item Multiple size- and value-based funds trade alongside S\&P 500
\end{itemize}

\vspace{0.3cm}
\begin{block}{Empirical Evidence}
\begin{itemize}
    \item Small stocks have equity premium $P$ highly correlated with large stock benchmark $P_0$
    \item High $R^2$ but $\beta$ slightly larger than 1
    \item Mid-cap: $\beta = 1.15$, Small-cap: $\beta = 1.27$
\end{itemize}
\end{block}
\end{frame}

\begin{frame}{CAPM-Based Model with Size Factor}
\begin{block}{Relative Size Definition}
$$C = \ln(S/S_0)$$
where $S$ = target portfolio market cap, $S_0$ = benchmark market cap
\end{block}

\begin{block}{Linear Model}
\begin{align*}
\alpha &= aC, \quad \beta = bC + 1\\
P &= AC + (1 + BC)P_0 + \delta
\end{align*}
\end{block}

\begin{itemize}
    \item $C = 0$: Target = benchmark ($\alpha = 0$, $\beta = 1$)
    \item Empirical: $a \approx 0$ but $b < 0$ (for small-cap and mid-cap)
\end{itemize}
\end{frame}

\section{Model with Volatility Normalization}

\begin{frame}{Price Returns Model}
\begin{block}{Complete Model}
$$R = aC + (1 + bC)R_0 + \varepsilon$$
where $R$ = price returns, $R_0$ = benchmark price returns
\end{block}

\begin{block}{Volatility Index (VIX)}
\begin{itemize}
    \item VIX measures market volatility expectations
    \item Dividing returns by VIX makes them independent and normal
    \item VIX modeled as autoregression of order 1 on logarithmic scale
\end{itemize}
\end{block}
\end{frame}

\begin{frame}{Discrete-Time Model Structure}
\begin{block}{Three Key Components}
\begin{enumerate}
    \item \textbf{Volatility}: VIX monthly average
    \item \textbf{Relative Size}: $C = \ln(S/S_0)$
    \item \textbf{CAPM}: Linear relationship with benchmark
\end{enumerate}
\end{block}

\begin{itemize}
    \item Uses real-world data from Kenneth French's data library
    \item Federal Reserve Economic Data (FRED) for VIX
    \item Proves long-term stability (ergodicity)
    \item Connects to Stochastic Portfolio Theory
\end{itemize}
\end{frame}

\section{Methodology}

\begin{frame}{Data and Estimation}
\begin{itemize}
    \item \textbf{Data Source}: Kenneth French Data Library + FRED
    \item \textbf{Portfolios}: Small-cap, mid-cap, large-cap indices
    \item \textbf{Benchmark}: S\&P 500 (large USA stocks)
    \item \textbf{Method}: Statistical fit of discrete-time model
    \item \textbf{Analysis}: Long-term stability conditions
\end{itemize}
\end{frame}

\section{Results}

\begin{frame}{Main Findings}
\begin{block}{Long-Term Stability}
\begin{itemize}
    \item Model exhibits ergodicity (long-term stability)
    \item Well-diversified portfolios stay together over time
    \item Capital distribution curve is stable
\end{itemize}
\end{block}

\begin{block}{Stochastic Portfolio Theory Connection}
\begin{itemize}
    \item Portfolios constructed as functions of market weights
    \item Small stocks have higher risk and return
    \item Distribution reduces to normal order statistics
\end{itemize}
\end{block}
\end{frame}

\begin{frame}{Empirical Results}
\small
\begin{table}
\centering
\begin{tabular}{lccc}
\hline
Portfolio & $\beta$ & $R^2$ & Observations \\
\hline
Small-cap & 1.27 & 0.89 & 240 \\
Mid-cap & 1.15 & 0.92 & 240 \\
Large-cap & 1.00 & 0.95 & 240 \\
\hline
\end{tabular}
\caption{CAPM Beta Estimates by Market Capitalization}
\end{table}
\end{frame}

\begin{frame}{Statistical Analysis}
\small
\begin{table}
\centering
\begin{tabular}{lcc}
\hline
Parameter & Estimate & Std. Error \\
\hline
$a$ (alpha) & 0.002 & 0.001 \\
$b$ (size factor) & -0.15 & 0.03 \\
Volatility coef. & 0.85 & 0.05 \\
\hline
\end{tabular}
\caption{Model Parameter Estimates}
\end{table}
\end{frame}

\section{Conclusion}

\begin{frame}{Key Contributions}
\begin{enumerate}
    \item \textbf{Unified Framework}: Combines volatility, size, and CAPM
    \item \textbf{Theoretical Rigor}: Proves long-term stability conditions
    \item \textbf{Empirical Validation}: Fits real market data successfully
    \item \textbf{Fills Research Gaps}: Extends previous work on CAPM with size factor
\end{enumerate}

\vspace{0.5cm}
\begin{block}{Practical Implications}
\begin{itemize}
    \item Better risk management for portfolio managers
    \item Improved understanding of size effect dynamics
    \item Framework for analyzing market stability
\end{itemize}
\end{block}
\end{frame}

\begin{frame}{Future Research Directions}
\begin{itemize}
    \item Extend to other factors (value, momentum, quality)
    \item Test on international markets
    \item Dynamic portfolio optimization applications
    \item High-frequency data analysis
\end{itemize}
\end{frame}

\begin{frame}
\centering
\vspace{2cm}
{\Huge Thank you for your attention!}\\
\vspace{1cm}
{\Large Questions?}
\end{frame}

\end{document}
